% !TeX root = ../papers/nc-p3.tex
\section{階数15の添字保持同型と二つの独立な不変量}
\label{sec:rank-fifteen-adjoint-invariants}
\begin{theorem}\label{thm:rank-fifteen-adjoint-invariants}
階数15の両表示の各々は、添字保持同型で割っても二次元の標準Koszul左右AS正則窓を与える。
具体的には通常窓から二つの有理不変量 $J_2,J_3$ が得られ、各表示で代数的に独立である。
この二値への写像は一般に有限対一である。
\end{theorem}

最初に同型の必要十分条件を記す。$M:V_1^*\otimes V_3^*\to V_0\otimes V_2$ を
反対切断とし、非零の右核と左核の生成元を $N,D$ とする。前定理によりいずれも
非退化な二文字形式である。$D$ で $V_2$ と $V_0^*$ を、$N$ で $V_3^*$ と $V_1$
を同一視すると、$M$ の核は恒等行列の直線、像はトレース零の行列空間となる。
従ってスカラーを除く線型同型
\[
 T:\mathfrak{sl}(V_1)\longrightarrow\mathfrak{sl}(V_0)
\]
を得る。第一位置を固定した二窓が添字保持同型である必要十分条件は、対応する二つの $T$ が
$T'=a\operatorname{Ad}(g_0)T\operatorname{Ad}(g_1)^{-1}$ を満たすこと
（$a\ne0$）である。実際、添字保持同型は最初の四文字空間に作用し、
前定理の $K_{0,4}=kW$ を保つ。核形式を正規化すると、二つの反対側の基底変更は
双対基底変更とスカラーに限られるのでこの式が必要となる。逆にこの式から $M$ の
スカラー付き基底同値を復元できる。対応する四階テンソルの巡回する隣接切断が
四つの関係空間を決めるため、同じ基底変更を周期的に繰り返せば添字保持同型を得る。

一つの四行四列基底変更だけで調べるには、$X,Y\in\mathfrak{sl}(V_1)$ に
\[
 b_T(X,Y)=\tfrac14\operatorname{tr}(TX\,TY),\qquad
 m_T(X,Y)=T^{-1}\bigl(TX\,TY-b_T(X,Y)I\bigr)
\]
を付す。$U=\operatorname{Ad}(g_1)$ と $a\ne0$ が存在して
\[
 b_{T'}(UX,UY)=a^2 b_T(X,Y),\qquad
 m_{T'}(UX,UY)=a\,U m_T(X,Y)
\]
が全基底対で成り立つことが同値条件である。必要性は代入で従う。逆に
$I\mapsto I,\ TX\mapsto T'(UX)/a$ は二式により積を保つ全行列代数の同型である。
行列代数の単位的自己同型は内部自己同型なので $g_0$ が得られる。
この条件は $\det(g_1)$ と $a$ を逆にした有限多項式系として検査できる。

\begin{proof}[定理の証明]
両トレース零空間の非退化形式 $\operatorname{tr}(XY)$ で随伴を取り、
$S=T^*T$ とする。基底同値で $S$ はスカラー倍と共役しか受けないので、
$\operatorname{tr}S\ne0$ では
\[
 J_m=\frac{\operatorname{tr}(S^m)}{(\operatorname{tr}S)^m},\qquad m=2,3
\]
が添字保持同型不変である。チャートで核直線の生成元は小行列式の比で取れるから、
これらは有理関数である。

$x=K+H$ と $c$ を相対曲面の座標として、$101$ で $x=4,c=3/4,L=5/4$、
$A=10,B=6$ のシートを取る。三平方根の初期値が非零であるため、
$x,c$ はこの点でエタール座標である。
\path{computations/verify_rank_fifteen_adjoint_invariants.py} は
$\mathbb Z/(101^2)$ 上で三根を一回のNewton操作で持ち上げる。
反対切断で単元の15次小行列式を選んで両核を計算し、元の行列との積が零であることを
確認する。基底には非対角行列単位と $E_{ii}-E_{44}$ の三本を取り、そのトレース形式
の行列を $G$ とすれば、$S$ の行列は $G^{-1}T^{\mathsf T}GT$ である。
始点での $(J_2,J_3)$ は両表示とも $(1418,4898)\bmod 101^2$ である。
$x$ のみ、または $c$ のみを101だけ動かし、同じシートで差を101で割ると、偏微分の二列は
$(22,88)^{\mathsf T},(23,66)^{\mathsf T}$ となる。分母は全て単元である。
行列式は $34\bmod101$ で、標数零でも二つの微分の外積は恒等的に零ではない。

従って両表示の各々で不変量写像の像の次元は2である。始域も既約な曲面なので、
一般ファイバーは零次元で、有限型性から有限である。
定理\ref{thm:rank-fifteen-vertical-curves}の有限集合を除いても像の次元は変わらない。
\end{proof}
